\documentclass[a4paper]{article}

\usepackage[spanish]{babel}
\usepackage{graphicx}
\usepackage{amsmath, amssymb}
\usepackage[margin=2cm]{geometry}
\usepackage{fancyhdr}
\usepackage{enumerate}
\usepackage[shortlabels]{enumitem}
\usepackage{parskip}
\usepackage[most]{tcolorbox}
\usepackage[hidelinks]{hyperref}
\usepackage{float}
\usepackage{booktabs}
\usepackage{mathrsfs}
\usepackage{tikz}

% cabecera
\pagestyle{fancy}
\fancyhead[l]{Astroestadística}
\fancyhead[c]{Parcial 1}
\fancyhead[r]{\today}
\fancyfoot[c]{\thepage}
\renewcommand{\headrulewidth}{0.2pt} % linea horizontal


\begin{document}


\section{Problema 1: Encuentro aleatorio}
Un evento astronómico raro ocurrirá en una ventana de observación específica de una hora. Dos observadores intentan registrarlo simultáneamente. El primero puede observar durante $10$ minutos desde que inicia la observación. El segundo puede observar $20$ minutos bajo las mismas condiciones. Suponga que sus tiempos de inicio de observación son independientes y uniformes.
\begin{enumerate}[a.]
	\item $(10\%)$ Sean  $T_1$ y $T_2$ los tiempos de inicio de los observadores $1$ y $2$ respectivamente. Escriba las funciones de distribución de probabilidad para $T_1$ y $T_2$ y su función de densidad conjunta.
	\item $(5\%)$ Haga una representación gráfica de $(t_1,t_2)$ para la que los observadores pueden estar observando simultáneamente.
	\item $(15\%)$ ¿Cuál es la probabilidad de que ambos observadores estén observando simultáneamente en algún instante?
\end{enumerate}

\subsection{Solución}
\begin{enumerate}[a.]
\item Tomamos el inicio de la ventana como $t=0$ y medimos los tiempos en minutos. Interpretamos que cada observador puede iniciar en cualquier instante de la hora de observación, incluso si parte de su intervalo queda fuera de la ventana. Por tanto:
$$T_1 \sim \mathrm{Unif}(0,60), \qquad T_2 \sim \mathrm{Unif}(0,60),$$
e independientes.

Las densidades marginales son
$$
f_{T_1}(t_1)=
\begin{cases}
\frac{1}{60}, & 0\le t_1\le 60,\\
0, & \text{en otro caso},
\end{cases}
\qquad
f_{T_2}(t_2)=
\begin{cases}
\frac{1}{60}, & 0\le t_2\le 60,\\
0, & \text{en otro caso}.
\end{cases}
$$

y las funciones de distribución acumulada son
$$
F_{T_1}(t_1)=
\begin{cases}
0, & t_1<0,\\
\frac{t_1}{60}, & 0\le t_1\le 60,\\
1, & t_1>60,
\end{cases}
\qquad
F_{T_2}(t_2)=
\begin{cases}
0, & t_2<0,\\
\frac{t_2}{60}, & 0\le t_2\le 60,\\
1, & t_2>60.
\end{cases}
$$

La densidad conjunta, por independencia, es uniforme sobre el cuadrado $[0,60]\times[0,60]$:
$$
f_{T_1,T_2}(t_1,t_2)=f_{T_1}(t_1)f_{T_2}(t_2)=
\begin{cases}
\frac{1}{3600}, & (t_1,t_2)\in[0,60]\times[0,60],\\
0, & \text{en otro caso}.
\end{cases}
$$

\item El observador 1 observa en el intervalo $[T_1,\,T_1+10]$ y el observador 2 en $[T_2,\,T_2+20]$ (minutos).
Hay simultaneidad en algún instante si y sólo si los intervalos se traslapan, es decir,
$$[T_1,T_1+10]\cap[T_2,T_2+20]\neq\varnothing.$$
Esto equivale a exigir simultáneamente
$$T_1 \le T_2+20 \qquad \text{y} \qquad T_2 \le T_1+10,$$
o, de manera equivalente,
$$T_2 \ge T_1-20, \qquad T_2 \le T_1+10,$$
junto con el soporte $0\le T_1\le 60$, $0\le T_2\le 60$.

\begin{center}
\begin{minipage}{0.95\linewidth}
\centering
\shorthandoff{<>}
\textbf{Si el observador 1 comienza primero: $T_1\le T_2$}

\begin{tikzpicture}[x=0.23cm,y=0.8cm,font=\small]
  \fill[green!18] (16,0.55) rectangle (20,2.45);
  \draw[dashed,gray] (16,0.55) -- (16,2.45);
  \draw[dashed,gray] (20,0.55) -- (20,2.45);
  \node[left] at (8,2) {Observador 1};
  \node[left] at (8,1) {Observador 2};
  \draw[blue,very thick,|-|] (10,2) -- (20,2)
    node[midway,below] {$10$ min};
  \node[above] at (10,2) {$T_1$};
  \node[above] at (20,2) {$T_1+10$};
  \draw[red,very thick,|-|] (16,1) -- (36,1)
    node[midway,above] {$20$ min};
  \node[below] at (16,1) {$T_2$};
  \node[below] at (36,1) {$T_2+20$};
  \draw[->] (8,0) -- (42,0) node[right] {tiempo};
  \node at (25,-0.6) {Hay traslape si $T_2\le T_1+10$.};
\end{tikzpicture}

\medskip
\textbf{Si el observador 2 comienza primero: $T_2\le T_1$}

\begin{tikzpicture}[x=0.23cm,y=0.8cm,font=\small]
  \fill[green!18] (24,0.55) rectangle (30,2.45);
  \draw[dashed,gray] (24,0.55) -- (24,2.45);
  \draw[dashed,gray] (30,0.55) -- (30,2.45);
  \node[left] at (8,2) {Observador 1};
  \node[left] at (8,1) {Observador 2};
  \draw[blue,very thick,|-|] (24,2) -- (34,2)
    node[midway,below] {$10$ min};
  \node[above] at (24,2) {$T_1$};
  \node[above] at (34,2) {$T_1+10$};
  \draw[red,very thick,|-|] (10,1) -- (30,1)
    node[midway,above] {$20$ min};
  \node[below] at (10,1) {$T_2$};
  \node[below] at (30,1) {$T_2+20$};
  \draw[->] (8,0) -- (42,0) node[right] {tiempo};
  \node at (25,-0.6) {Hay traslape si $T_1\le T_2+20$, es decir, $T_2\ge T_1-20$.};
\end{tikzpicture}
\shorthandon{<>}

\small La franja verde indica el intervalo común en estos dos ejemplos de traslape. Los segmentos azules duran $10$ minutos y los rojos $20$ minutos; las posiciones de inicio son ilustrativas.
\end{minipage}
\end{center}

Para comprobar el sentido de estos límites, distinguimos quién inicia primero. Si $T_2\ge T_1$, el observador 1 debe seguir observando cuando comienza el 2, por lo que $T_2-T_1\le 10$. Si $T_2\le T_1$, el observador 2 debe seguir observando cuando comienza el 1, por lo que $T_1-T_2\le 20$. Así, el límite superior es $t_2=t_1+10$ y el inferior es $t_2=t_1-20$; las duraciones no se intercambian.

Además, cuando hay traslape, ambos están observando en el instante $\max\{T_1,T_2\}$, que pertenece a $[0,60]$. Por ello, permitir que una sesión termine después del minuto $60$ no introduce coincidencias que ocurran exclusivamente fuera de la ventana.


Gráficamente, en el plano $(t_1,t_2)$, la región de simultaneidad es la intersección del cuadrado $[0,60]\times[0,60]$ con la franja delimitada por las rectas $t_2=t_1-20$ y $t_2=t_1+10$:

Los vértices de esta región, recorridos en orden, son
$$
(0,0),\ (20,0),\ (60,40),\ (60,60),\ (50,60),\ (0,10).
$$

\begin{center}
\shorthandoff{<>}
\begin{tikzpicture}[scale=0.09]
  % axes
  \draw[->] (-3,0) -- (67,0) node[below right] {$t_1$ (min)};
  \draw[->] (0,-3) -- (0,67) node[above left] {$t_2$ (min)};

  % support square
  \draw[thick] (0,0) rectangle (60,60);

  % boundary lines of the band
  \draw[blue, thick] (20,0) -- (60,40)
    node[pos=0.92, below right, fill=white, inner sep=1pt] {$t_2=t_1-20$};
  \draw[red, thick] (0,10) -- (50,60)
    node[pos=0.78, above left, fill=white, inner sep=1pt] {$t_2=t_1+10$};

  % feasible region: (0,0)-(20,0)-(60,40)-(60,60)-(50,60)-(0,10)
  \fill[gray!25] (0,0) -- (20,0) -- (60,40) -- (60,60) -- (50,60) -- (0,10) -- cycle;

  % vertices
  \foreach \x/\y in {0/0,20/0,60/40,60/60,50/60,0/10}{
    \fill (\x,\y) circle (1.15);
  }

  % vertex labels (offset to avoid overlap)
  \node[below left]  at (0,0)   {$(0,0)$};
  \node[below]       at (20,0)  {$(20,0)$};
  \node[right]       at (60,40) {$(60,40)$};
  \node[above right] at (60,60) {$(60,60)$};
  \node[above]       at (50,60) {$(50,60)$};
  \node[left]        at (0,10)  {$(0,10)$};

  % region label
  \node[fill=white, inner sep=1.2pt] at (26,31) {Región de simultaneidad};
\end{tikzpicture}
\shorthandon{<>}
\end{center}

\item \textbf{Justificación de la equivalencia geométrica.}
Como los tiempos de inicio son independientes y uniformes en $[0,60]$, su densidad conjunta es constante:
$$
f_{T_1,T_2}(t_1,t_2)=\frac{1}{3600},
\qquad (t_1,t_2)\in S=[0,60]\times[0,60].
$$
Por definición, la probabilidad de que el par $(T_1,T_2)$ pertenezca a una región medible $D\subseteq S$ se obtiene integrando la densidad conjunta sobre ella. Al ser constante, puede extraerse de la integral:
$$
P\bigl((T_1,T_2)\in D\bigr)
=\iint_D f_{T_1,T_2}(t_1,t_2)\,dt_1\,dt_2
=\frac{1}{3600}\iint_D 1\,dt_1\,dt_2
=\frac{\text{Área}(D)}{3600}.
$$
Así, la probabilidad es el cociente entre el área de la región favorable y el área total del cuadrado, $\text{Área}(S)=60\cdot60=3600$. Esta equivalencia se debe a la uniformidad de la densidad conjunta; no es válida en general para una densidad no constante.

\textbf{Región de simultaneidad.}
Los intervalos de observación se intersectan si y sólo si
$$
T_1\le T_2+20
\quad\text{y}\quad
T_2\le T_1+10,
$$
es decir, si $T_1-20\le T_2\le T_1+10$. Por tanto, la región favorable es
$$
D=\bigl\{(t_1,t_2)\in S: t_1-20\le t_2\le t_1+10\bigr\}.
$$

\textbf{Cálculo del área mediante el complemento.}
En lugar de dividir la región sombreada en varias partes, restamos del cuadrado los dos triángulos donde no hay simultaneidad:
\begin{itemize}
\item Si $t_2>t_1+10$, el segundo observador comienza después de que el primero haya terminado. La recta $t_2=t_1+10$ corta al cuadrado en $(0,10)$ y $(50,60)$. La región excluida es un triángulo rectángulo con vértices $(0,10)$, $(0,60)$ y $(50,60)$, cuyos catetos miden $50$ minutos. Su área es
$$A_1=\frac{1}{2}(60-10)^2=1250.$$
\item Si $t_2<t_1-20$, el primero comienza después de que el segundo haya terminado. La recta $t_2=t_1-20$ corta al cuadrado en $(20,0)$ y $(60,40)$. La región excluida es un triángulo rectángulo con vértices $(20,0)$, $(60,0)$ y $(60,40)$, cuyos catetos miden $40$ minutos. Su área es
$$A_2=\frac{1}{2}(60-20)^2=800.$$
\end{itemize}
Estas dos regiones son disjuntas y abarcan todos los casos sin traslape. Las rectas frontera tienen área, y por tanto probabilidad, cero, de modo que incluir o excluir los casos en que los intervalos sólo comparten un extremo no altera el resultado. En consecuencia,
$$
\text{Área}(D)=3600-1250-800=1550,
$$
y la probabilidad buscada es
$$
\boxed{P(\text{simultaneidad})=\frac{1550}{3600}=\frac{31}{72}\approx 0.4306.}
$$

\textbf{Comprobación por integración directa.}
Para cada $t_1\in[0,60]$, la sección vertical de la región correcta satisface
$$
\max\{0,t_1-20\}\le t_2\le\min\{60,t_1+10\}.
$$
El límite inferior cambia en $t_1=20$ y el superior en $t_1=50$. Por tanto,
\begin{align*}
\text{Área}(D)
&=\int_0^{20}(t_1+10)\,dt_1
 +\int_{20}^{50}30\,dt_1
 +\int_{50}^{60}(80-t_1)\,dt_1\\
&=400+900+250=1550,
\end{align*}
lo que confirma tanto la región representada como la probabilidad obtenida por el complemento.

\textit{Interpretación:} bajo el supuesto de que ambos pueden iniciar en cualquier instante de la hora, la probabilidad de que coincidan observando en algún instante es aproximadamente $43.06\%$. Se calcula la probabilidad de traslape de sus intervalos, no la de registrar ambos el evento astronómico, para la cual haría falta modelar también el instante en que ocurre.
\end{enumerate}

\section{Problema 2: Estructura galáctica}
En un modelo simplificado del disco de una galaxia, la distribución vertical de estrellas respecto al plano galáctico se describe mediante la siguiente función de distribución de probabilidad:

$$f_Z(z) = \frac{1}{2h}e^{-|z|/h}$$

con $z\in\mathbb{R}$ y $h>0$ es la escala de altura del disco.

Suponga que la distancia radial efectiva $R$ de una estrella al centro galáctico depende de su altura $Z$ según la relación

$$R=R_0 + \alpha Z^2$$

donde $R_0>0$ y $\alpha >0$ son constantes.

\begin{enumerate}[a.]
	\item $(15\%)$ Determine la función de distribución de probabilidad de $R$, $g(R)$.
	\item $(10\%)$ Verifique explícitamente que la función de densidad de probabilidad obtenida está correctamente normalizada.
	\item Calcule los valores esperados $E[Z]$ y $E[R]$
	
	\textit{Ayuda:} $$\int_{0}^{\infty}u^ne^{-u}du=n!,\quad n\in\mathbb{N}$$
\end{enumerate}

\subsection{Solución}

\begin{enumerate}[a.]
\item \textbf{Valores de $Z$ que producen un mismo valor de $R$.}
La variable $Z$ tiene densidad
$$
f_Z(z)=\frac{1}{2h}e^{-|z|/h},\qquad z\in\mathbb{R},\quad h>0,
$$
y se transforma mediante $R=R_0+\alpha Z^2$, con $R_0>0$ y $\alpha>0$. Como $Z^2\ge0$, necesariamente $R\ge R_0$. Por tanto, el soporte de $R$ es $[R_0,\infty)$ y su densidad es cero para $r<R_0$.

Para cada $r>R_0$, despejamos $z$ en la relación dada:
$$
r=R_0+\alpha z^2
\quad\Longleftrightarrow\quad
z^2=\frac{r-R_0}{\alpha}.
$$
Existen dos valores de $z$ que producen el mismo $r$, uno positivo y otro negativo:
$$
z_1(r)=\sqrt{\frac{r-R_0}{\alpha}},
\qquad
z_2(r)=-\sqrt{\frac{r-R_0}{\alpha}}.
$$

\begin{figure}[H]
\centering
\shorthandoff{<>}
\begin{minipage}[t]{0.48\linewidth}
\centering
\textbf{(a) Distribución de alturas}

\begin{tikzpicture}[x=0.83cm,y=6.2cm,font=\small]
  \fill[blue!10] (-3,0) -- plot[domain=-3:0,samples=60,variable=\heightvalue]
    ({\heightvalue},{0.5*exp(\heightvalue)})
    -- plot[domain=0:3,samples=60,variable=\heightvalue]
    ({\heightvalue},{0.5*exp(-\heightvalue)}) -- (3,0) -- cycle;
  \draw[->] (-3.4,0) -- (3.5,0) node[right] {$s=z/h$};
  \draw[->] (0,0) -- (0,0.64) node[above] {$h f_Z(hs)$};
  \draw[blue!75!black,very thick] plot[domain=-3:0,samples=60,variable=\heightvalue]
    ({\heightvalue},{0.5*exp(\heightvalue)});
  \draw[blue!75!black,very thick] plot[domain=0:3,samples=60,variable=\heightvalue]
    ({\heightvalue},{0.5*exp(-\heightvalue)});
  \draw[dashed,gray] (-1,0) -- (-1,0.18394) -- (1,0.18394) -- (1,0);
  \fill[red!75!black] (-1,0.18394) circle (2pt);
  \fill[red!75!black] (1,0.18394) circle (2pt);
  \node[above left] at (0,0.5) {$1/2$};
  \foreach \tick in {-3,-2,-1,1,2,3}{
    \draw (\tick,0) -- (\tick,-0.008) node[below] {$\tick$};
  }
  \node[below left] at (0,0) {$0$};
  \node[align=center] at (0,-0.14) {$f_Z(-h)=f_Z(h)$\\$E[Z]=0$ por simetría};
\end{tikzpicture}
\end{minipage}\hfill
\begin{minipage}[t]{0.48\linewidth}
\centering
\textbf{(b) Transformación cuadrática}

\begin{tikzpicture}[x=1.05cm,y=0.95cm,font=\small]
  \draw[->] (-2.35,0) -- (2.5,0) node[right] {$s=z/h$};
  \draw[->] (0,0) -- (0,4.2) node[above] {$x=\dfrac{r-R_0}{\alpha h^2}$};
  \draw[blue!75!black,very thick] plot[domain=-1.9:1.9,samples=100,variable=\heightvalue]
    ({\heightvalue},{\heightvalue*\heightvalue});
  \draw[red!75!black,dashed] (-2.1,1) -- (2.1,1) node[right] {$x=1$};
  \draw[dashed,gray] (-1,0) -- (-1,1);
  \draw[dashed,gray] (1,0) -- (1,1);
  \fill[red!75!black] (-1,1) circle (2pt);
  \fill[red!75!black] (1,1) circle (2pt);
  \foreach \tick in {-2,-1,1,2}{
    \draw (\tick,0) -- (\tick,-0.05) node[below] {$\tick$};
  }
  \node[below left] at (0,0) {$0$};
  \node at (0.9,3.4) {$x=s^2$};
  \node[align=center] at (0,-0.9) {$z=-h$ y $z=h$ producen\\el mismo $r=R_0+\alpha h^2$};
\end{tikzpicture}
\end{minipage}
\shorthandon{<>}
\caption{La distribución de alturas es simétrica respecto al plano galáctico. Dos alturas opuestas producen el mismo radio efectivo; los puntos rojos ilustran el caso $z=\pm h$. Por ello, ambas alturas deben contribuir al cálculo de $g(r)$. Los ejes están expresados en variables adimensionales.}
\label{fig:altura-transformacion}
\end{figure}

\textbf{Ecuación para construir la densidad $g(r)$.}
Como tanto $z_1(r)$ como $z_2(r)$ producen el mismo valor de $r$, debemos sumar sus contribuciones. La fórmula de transformación de densidades es
$$
\boxed{
g(r)=f_Z\bigl(z_1(r)\bigr)\left|\frac{dz_1}{dr}\right|
     +f_Z\bigl(z_2(r)\bigr)\left|\frac{dz_2}{dr}\right|,
\qquad r>R_0.}
$$
Cada factor $|dz_i/dr|$ es el Jacobiano correspondiente: convierte una pequeña longitud en la variable $r$ en la longitud correspondiente en $z$. Se toma su valor absoluto porque las longitudes y las contribuciones a la probabilidad son no negativas. A continuación calculamos los factores de esta ecuación.

\textbf{Cálculo de los Jacobianos.}
Derivamos primero la relación $r=R_0+\alpha z^2$:
$$
\frac{dr}{dz}=2\alpha z
\quad\Longrightarrow\quad
\left|\frac{dz}{dr}\right|=\frac{1}{2\alpha|z|},
\qquad z\ne0.
$$
Como
$$
|z_1(r)|=|z_2(r)|=\sqrt{\frac{r-R_0}{\alpha}},
$$
los dos Jacobianos son iguales:
$$
\left|\frac{dz_1}{dr}\right|
=\left|\frac{dz_2}{dr}\right|
=\frac{1}{2\alpha\sqrt{(r-R_0)/\alpha}}
=\frac{1}{2\sqrt{\alpha(r-R_0)}},
\qquad r>R_0.
$$

\textbf{Sustitución en la ecuación de $g(r)$.}
La densidad de $Z$ depende de $|z|$, de modo que
$$
f_Z\bigl(z_1(r)\bigr)=f_Z\bigl(z_2(r)\bigr)
=\frac{1}{2h}\exp\!\left(-\frac{1}{h}\sqrt{\frac{r-R_0}{\alpha}}\right).
$$
Por tanto, los dos términos de la ecuación de $g(r)$ son iguales. Al sustituir las densidades y los Jacobianos, obtenemos
\begin{align*}
g(r)
&=2\left[\frac{1}{2h}\exp\!\left(-\frac{1}{h}\sqrt{\frac{r-R_0}{\alpha}}\right)\right]
  \frac{1}{2\sqrt{\alpha(r-R_0)}}\\
&=\frac{1}{2h\sqrt{\alpha(r-R_0)}}
  \exp\!\left(-\frac{1}{h}\sqrt{\frac{r-R_0}{\alpha}}\right),
\qquad r>R_0.
\end{align*}
En consecuencia, la densidad buscada puede escribirse como
$$
\boxed{
g(r)=
\begin{cases}
\dfrac{1}{2h\sqrt{\alpha(r-R_0)}}
\exp\!\left(-\dfrac{1}{h}\sqrt{\dfrac{r-R_0}{\alpha}}\right), & r>R_0,\\[4mm]
0, & r\le R_0.
\end{cases}}
$$

\textit{Observación sobre $r=R_0$:} este valor corresponde únicamente a $z=0$, donde no se puede aplicar la expresión anterior del Jacobiano. Aunque la densidad crece sin límite al acercarse a $R_0$ por la derecha, su integral es finita, como se comprobará en el siguiente inciso. Además, $P(R=R_0)=P(Z=0)=0$, porque $Z$ es continua. Por ello, asignar $g(R_0)=0$ no altera ninguna probabilidad.

\begin{figure}[H]
\centering
\shorthandoff{<>}
\begin{tikzpicture}[x=1.1cm,y=1.05cm,font=\small]
  \fill[green!15] (0.0144,0) --
    plot[domain=0.12:2.828427,samples=200,variable=\rootvalue]
    ({\rootvalue*\rootvalue},{exp(-\rootvalue)/(2*\rootvalue)})
    -- (8,0) -- cycle;
  \draw[->] (-0.3,0) -- (8.6,0) node[below] {$x$};
  \draw[->] (0,0) -- (0,4.3) node[above] {$p_X(x)$};
  \draw[blue!75!black,very thick,<-]
    plot[domain=0.12:2.828427,samples=200,variable=\rootvalue]
    ({\rootvalue*\rootvalue},{exp(-\rootvalue)/(2*\rootvalue)});
  \foreach \tick in {1,2,3,4,5,6,7,8}{
    \draw (\tick,0) -- (\tick,-0.06) node[below] {$\tick$};
  }
  \foreach \tick in {1,2,3,4}{
    \draw (0,\tick) -- (-0.06,\tick) node[left] {$\tick$};
  }
  \node[below left] at (0,0) {$0$};
  \draw[orange!80!black,dashed,thick] (2,0) -- (2,1.1)
    node[above] {$E[X]=2$};
  \node[align=center] at (5.1,3.3)
    {$\displaystyle p_X(x)=\frac{e^{-\sqrt{x}}}{2\sqrt{x}},\quad x>0$};
  \node[align=center] at (5.1,2.3)
    {$\displaystyle\int_0^\infty p_X(x)\,dx=1$\\[6pt]Área total bajo la curva};
  \node[align=center,text=blue!75!black] at (1.8,3.6)
    {Crece sin límite\\cuando $x\to0^+$};
  \draw[->,blue!75!black] (1,3.5) -- (0.08,3.5);
\end{tikzpicture}
\shorthandon{<>}

\smallskip
\(\displaystyle X=\frac{R-R_0}{\alpha h^2},\qquad
p_X(x)=\alpha h^2\,g(R_0+\alpha h^2x).\)
\caption{Densidad del radio efectivo en escala adimensional. La divergencia en el extremo izquierdo es integrable y no representa una probabilidad puntual en $R_0$. La curva continúa para $x>8$; su área total es uno. La línea discontinua señala la media $E[X]=2$, equivalente a $E[R]=R_0+2\alpha h^2$.}
\label{fig:densidad-radio}
\end{figure}

\item Verificamos la normalización:
$$
\int_{-\infty}^{\infty} g(r)\,dr=\int_{R_0}^{\infty}\frac{1}{2h\sqrt{\alpha(r-R_0)}}
\exp\!\left(-\frac{1}{h}\sqrt{\frac{r-R_0}{\alpha}}\right)\,dr.
$$
Hacemos el cambio
$$
u=\sqrt{\frac{r-R_0}{\alpha}}\quad\Longrightarrow\quad r=R_0+\alpha u^2,\ \ dr=2\alpha u\,du,
$$
y además $\sqrt{\alpha(r-R_0)}=\sqrt{\alpha\cdot \alpha u^2}=\alpha u$ (para $u\ge 0$). Entonces
$$
\int_{R_0}^{\infty} g(r)\,dr
=\int_{0}^{\infty}\frac{1}{2h(\alpha u)}e^{-u/h}\,(2\alpha u)\,du
=\int_0^\infty \frac{1}{h}e^{-u/h}\,du
=\left[-e^{-u/h}\right]_{0}^{\infty}=1.
$$
Por tanto, $g(r)$ está correctamente normalizada.

\item Primero,
$$
E[Z]=\int_{-\infty}^{\infty} z\,\frac{1}{2h}e^{-|z|/h}\,dz=0,
$$
porque la densidad es par y la función $z$ es impar (distribución simétrica alrededor de $0$).

Para $E[R]$, usamos $R=R_0+\alpha Z^2$:
$$
E[R]=E[R_0+\alpha Z^2]=R_0+\alpha E[Z^2].
$$
Calculamos $E[Z^2]$:
$$
E[Z^2]=\int_{-\infty}^{\infty} z^2\,\frac{1}{2h}e^{-|z|/h}\,dz
=2\int_{0}^{\infty} z^2\,\frac{1}{2h}e^{-z/h}\,dz
=\int_0^\infty \frac{1}{h}z^2 e^{-z/h}\,dz.
$$
Con el cambio $u=z/h$ (i.e. $z=hu$, $dz=h\,du$),
$$
E[Z^2]=\int_0^\infty \frac{1}{h}(h^2u^2)e^{-u}(h\,du)
=h^2\int_0^\infty u^2 e^{-u}\,du.
$$
Usando la ayuda $\int_0^\infty u^n e^{-u}\,du=n!$ con $n=2$:
$$
\int_0^\infty u^2 e^{-u}\,du=2!=2,
$$
por tanto
$$
E[Z^2]=2h^2
\quad\Longrightarrow\quad
E[R]=R_0+\alpha(2h^2)=R_0+2\alpha h^2.
$$

\textit{Interpretación:} el disco es simétrico respecto al plano ($E[Z]=0$), pero la relación cuadrática $R=R_0+\alpha Z^2$ hace que, en promedio, la distancia radial efectiva aumente en $2\alpha h^2$ respecto a $R_0$, controlado por el espesor vertical (vía $h$) y el acoplamiento $\alpha$.
\end{enumerate}

\section{Problema 3: Detección de exoplanetas y actualización bayesiana}
En un programa de búsqueda de exoplanetas, se estima que el $15\%$ de las estrellas observadas albergan un planeta detectable mediante técnicas actuales, mientras que el resto no presenta señales planetarias. Considere que mediante el método de velocidad radial, se detecta correctamente un planeta con probabilidad de $0.9$ y se produce una detección falsa con probabilidad $0.05$ cuando no hay planeta.

\begin{enumerate}[a.]
	\item $(10\%)$ Si una estrella es clasificada como \textit{con planeta} por este método, ¿cuál es la probabilidad de que en realidad no tenga planeta?
	\item Para confirmar la detección, se utiliza un segundo método independiente (por ejemplo, tránsitos), con las mismas características estadísticas. Si este otro método también indica la presencia de un planeta, ¿Cuál es la probabilidad de que la detección sea falsa?
	\item Muestre que bajo independencia condicional de los métodos, dado el estado real del sistema (con o sin planeta), el resultado de aplicar los métodos secuenciales es equivalente a considerarlos simultáneamente.
\end{enumerate}

\subsection{Solución}
\begin{enumerate}[a.]
\item Sea $P$ el evento ``la estrella tiene planeta detectable'' y $\bar P$ su complemento. Sea $D_1$ el evento ``el método 1 clasifica como con planeta''.

Datos:
$$P(P)=0.15,\qquad P(\bar P)=0.85,$$
$$P(D_1\mid P)=0.9,\qquad P(D_1\mid \bar P)=0.05.$$

Se pide $P(\bar P\mid D_1)$. Por Bayes:
$$
P(\bar P\mid D_1)=\frac{P(D_1\mid \bar P)P(\bar P)}{P(D_1)}.
$$
y por probabilidad total,
$$
P(D_1)=P(D_1\mid P)P(P)+P(D_1\mid \bar P)P(\bar P)
=0.9(0.15)+0.05(0.85)=0.135+0.0425=0.1775.
$$
Entonces
$$
P(\bar P\mid D_1)=\frac{0.05(0.85)}{0.1775}=\frac{0.0425}{0.1775}
=\frac{85}{355}=\frac{17}{71}\approx 0.2394.
$$

\textit{Interpretación:} aun con una eficiencia alta ($0.9$) el porcentaje de falsos positivos a posteriori no es despreciable porque la prevalencia de planetas ($15\%$) es relativamente baja.

\begin{figure}[H]
\centering
\shorthandoff{<>}
\begin{tikzpicture}[
  font=\small,
  groupbox/.style={draw,rounded corners=3pt,align=center,minimum height=0.9cm,text width=2.6cm,inner sep=5pt},
  connection/.style={->,thick}
]
  \node[groupbox,fill=gray!10,text width=3.6cm] (total) at (0,0)
    {\textbf{8000 estrellas}\\Antes de observar};
  \node[groupbox,draw=blue!65!black,fill=blue!8] (planet) at (-3.2,-1.7)
    {\textbf{1200 con planeta}\\$15\%$ del total};
  \node[groupbox,draw=orange!80!black,fill=orange!12] (noplanet) at (3.2,-1.7)
    {\textbf{6800 sin planeta}\\$85\%$ del total};
  \draw[connection] (total) -- (planet);
  \draw[connection] (total) -- (noplanet);

  \node[groupbox,draw=blue!65!black,fill=blue!15] (truepositive) at (-4.8,-3.7)
    {\textbf{1080 positivos}\\Con planeta};
  \node[groupbox,draw=gray,fill=gray!7] (falsenegative) at (-1.6,-3.7)
    {120 negativos\\Con planeta};
  \node[groupbox,draw=orange!80!black,fill=orange!25] (falsepositive) at (1.6,-3.7)
    {\textbf{340 positivos}\\Sin planeta};
  \node[groupbox,draw=gray,fill=gray!7] (truenegative) at (4.8,-3.7)
    {6460 negativos\\Sin planeta};
  \draw[connection,blue!65!black] (planet) -- (truepositive)
    node[midway,left,fill=white,inner sep=2pt] {$90\%$};
  \draw[connection,gray] (planet) -- (falsenegative)
    node[midway,right,fill=white,inner sep=2pt] {$10\%$};
  \draw[connection,orange!80!black] (noplanet) -- (falsepositive)
    node[midway,left,fill=white,inner sep=2pt] {$5\%$};
  \draw[connection,gray] (noplanet) -- (truenegative)
    node[midway,right,fill=white,inner sep=2pt] {$95\%$};

  \node[draw,rounded corners=3pt,fill=green!7,align=center,text width=12cm,inner sep=9pt] at (0,-5.6)
    {Al conocer un resultado positivo, sólo cuentan las dos cajas coloreadas de la última fila.\\[5pt]
    $\displaystyle P(\bar P\mid D_1)
      =\frac{\textcolor{orange!80!black}{340}}
      {\textcolor{blue!65!black}{1080}+\textcolor{orange!80!black}{340}}
      =\frac{17}{71}\approx23.94\%$};
\end{tikzpicture}
\shorthandon{<>}
\caption{Bayes mediante frecuencias esperadas en una población ilustrativa de 8000 estrellas, no mediante datos observados. El $5\%$ se aplica a las 6800 estrellas sin planeta; el $23.94\%$ se calcula entre los 1420 resultados positivos. Son probabilidades con distintos denominadores.}
\label{fig:bayes-frecuencias}
\end{figure}

\item Ahora se aplica un segundo método independiente con las mismas características. Sea $D_2$ el evento ``el método 2 clasifica como con planeta''. Se asume independencia condicional dado el estado real:
$$
P(D_1\cap D_2\mid P)=P(D_1\mid P)P(D_2\mid P),\qquad
P(D_1\cap D_2\mid \bar P)=P(D_1\mid \bar P)P(D_2\mid \bar P).
$$
Como tiene las mismas características:
$$P(D_2\mid P)=0.9,\qquad P(D_2\mid \bar P)=0.05.$$

Se pide $P(\bar P\mid D_1\cap D_2)$. Por Bayes:
$$
P(\bar P\mid D_1\cap D_2)=\frac{P(D_1\cap D_2\mid \bar P)P(\bar P)}{P(D_1\cap D_2)}.
$$
Calculamos:
$$
P(D_1\cap D_2\mid P)=0.9\cdot 0.9=0.81,
\qquad
P(D_1\cap D_2\mid \bar P)=0.05\cdot 0.05=0.0025.
$$
Luego,
$$
P(D_1\cap D_2)=0.81(0.15)+0.0025(0.85)=0.1215+0.002125=0.123625.
$$
Por tanto,
$$
P(\bar P\mid D_1\cap D_2)=\frac{0.0025(0.85)}{0.123625}
=\frac{0.002125}{0.123625}
=\frac{17}{989}\approx 0.01719.
$$

También podemos obtener este resultado de manera secuencial, aprovechando la información del primer positivo. En una actualización bayesiana, las probabilidades posteriores obtenidas al incorporar una observación se utilizan como punto de partida para incorporar la siguiente. En este caso, el estado que queremos conocer sigue siendo el mismo: si la estrella tiene o no un planeta detectable.

Del inciso a sabemos que
$$
P(\bar P\mid D_1)=\frac{17}{71},
\qquad
P(P\mid D_1)=1-\frac{17}{71}=\frac{54}{71}.
$$
Estas probabilidades ya incluyen la información de $D_1$ y sustituyen a las probabilidades iniciales $0.85$ y $0.15$ al actualizar con el segundo positivo. La independencia condicional permite escribir
$$
P(D_2\mid \bar P,D_1)=P(D_2\mid\bar P)=0.05,
\qquad
P(D_2\mid P,D_1)=P(D_2\mid P)=0.9.
$$
Es decir, una vez fijado si la estrella tiene o no planeta, conocer el primer resultado no modifica las probabilidades de respuesta del segundo método.

Aplicando Bayes dentro del conjunto de estrellas con un primer resultado positivo, obtenemos
$$
P(\bar P\mid D_1\cap D_2)
=\frac{P(D_2\mid\bar P)P(\bar P\mid D_1)}
{P(D_2\mid\bar P)P(\bar P\mid D_1)+P(D_2\mid P)P(P\mid D_1)}.
$$
El denominador es $P(D_2\mid D_1)$: la probabilidad de obtener el segundo positivo cuando ya se conoce el primero. Al sustituir los valores anteriores,
$$
P(\bar P\mid D_1\cap D_2)
=\frac{0.05\left(\frac{17}{71}\right)}
{0.05\left(\frac{17}{71}\right)+0.9\left(\frac{54}{71}\right)}.
$$
Multiplicando el numerador y el denominador por $71$ y simplificando, resulta
\begin{align*}
P(\bar P\mid D_1\cap D_2)
&=\frac{0.05(17)}{0.05(17)+0.9(54)}
 =\frac{0.85}{0.85+48.6}\\
&=\frac{0.85}{49.45}
 =\frac{85}{4945}
 =\frac{17}{989}\approx0.01719.
\end{align*}
Así, actualizar con el primer positivo y luego con el segundo conduce al mismo resultado que considerar ambos positivos conjuntamente: la probabilidad de que la estrella no tenga planeta después de las dos detecciones es aproximadamente $1.72\%$.

\textit{Interpretación:} exigir confirmación por dos métodos reduce drásticamente la probabilidad de falso positivo (de $\approx 0.239$ a $\approx 0.017$), porque la probabilidad de doble falso positivo bajo $\bar P$ es $0.05^2$.

\begin{figure}[H]
\centering
\shorthandoff{<>}
\begin{tikzpicture}[font=\small]
  \node[anchor=west,font=\bfseries] at (0,0.7)
    {La segunda prueba se aplica a los 1420 primeros positivos};
  \node[anchor=west,text=blue!65!black] at (0,0) {1080 con planeta};
  \draw[->,thick,blue!65!black] (3.1,0) -- (6.1,0)
    node[midway,above] {$\times\,0.9$};
  \node[anchor=west,text=blue!65!black] at (6.3,0) {972 vuelven a dar positivo};
  \node[anchor=west,text=orange!80!black] at (0,-0.9) {340 sin planeta};
  \draw[->,thick,orange!80!black] (3.1,-0.9) -- (6.1,-0.9)
    node[midway,above] {$\times\,0.05$};
  \node[anchor=west,text=orange!80!black] at (6.3,-0.9) {17 vuelven a dar positivo};

  \node[anchor=west,font=\bfseries] at (0,-2) {Después del primer positivo: 1420 estrellas};
  \fill[blue!20] (0,-3.1) rectangle ({10*1080/1420},-2.3);
  \fill[orange!40] ({10*1080/1420},-3.1) rectangle (10,-2.3);
  \draw[thick] (0,-3.1) rectangle (10,-2.3);
  \node[align=center] at ({5*1080/1420},-2.7) {1080 con planeta\\$76.06\%$};
  \node[align=center] at ({5+5*1080/1420},-2.7) {340 sin planeta\\$23.94\%$};

  \node[anchor=west,font=\bfseries] at (0,-4.2) {Después de dos positivos: 989 estrellas};
  \fill[blue!20] (0,-5.3) rectangle ({10*972/989},-4.5);
  \fill[orange!40] ({10*972/989},-5.3) rectangle (10,-4.5);
  \draw[thick] (0,-5.3) rectangle (10,-4.5);
  \node[align=center] at ({5*972/989},-4.9) {972 con planeta\\$98.28\%$};
  \draw[->,orange!80!black] (10.4,-4.9) -- ({5+5*972/989},-4.9);
  \node[anchor=west,align=left,text=orange!80!black] at (10.45,-4.9)
    {17 sin planeta\\$1.72\%$};

  \node[align=center] at (5,-6.25)
    {$\displaystyle P(\bar P\mid D_1\cap D_2)=\frac{17}{972+17}=\frac{17}{989}$\\[6pt]
    La fracción naranja disminuye: $23.94\%\ \longrightarrow\ 1.72\%$.};
\end{tikzpicture}
\shorthandon{<>}
\caption{Confirmación secuencial bajo independencia condicional dado el estado real de la estrella. Se conservan el $90\%$ de los positivos con planeta y el $5\%$ de los positivos sin planeta. Las frecuencias son esperadas para la misma población ilustrativa de 8000 estrellas. Cada barra representa el $100\%$ de su propio grupo: sus anchos comparan proporciones, no cantidades de estrellas.}
\label{fig:bayes-confirmacion}
\end{figure}

\item (Equivalencia secuencial vs simultánea bajo independencia condicional)

\textbf{Enfoque simultáneo.} Considerar los resultados conjuntos $(D_1,D_2)$ conduce a la posterior:
$$
P(P\mid D_1\cap D_2)=\frac{P(D_1\cap D_2\mid P)P(P)}{P(D_1\cap D_2\mid P)P(P)+P(D_1\cap D_2\mid \bar P)P(\bar P)}.
$$
Bajo independencia condicional,
$$
P(D_1\cap D_2\mid P)=P(D_1\mid P)P(D_2\mid P),\qquad
P(D_1\cap D_2\mid \bar P)=P(D_1\mid \bar P)P(D_2\mid \bar P).
$$

\textbf{Enfoque secuencial.} Primero se actualiza con $D_1$:
$$
P(P\mid D_1)=\frac{P(D_1\mid P)P(P)}{P(D_1\mid P)P(P)+P(D_1\mid \bar P)P(\bar P)}.
$$
Luego, usando Bayes otra vez al incorporar $D_2$:
$$
P(P\mid D_1\cap D_2)=\frac{P(D_2\mid P, D_1)\,P(P\mid D_1)}{P(D_2\mid P, D_1)\,P(P\mid D_1)+P(D_2\mid \bar P, D_1)\,P(\bar P\mid D_1)}.
$$
La independencia condicional dado el estado real implica
$$
P(D_2\mid P, D_1)=P(D_2\mid P),\qquad P(D_2\mid \bar P, D_1)=P(D_2\mid \bar P).
$$
Sustituyendo y reemplazando $P(P\mid D_1)$ y $P(\bar P\mid D_1)$ por sus expresiones de Bayes, se obtiene (mostrando el paso clave de simplificación):
$$
P(P\mid D_1\cap D_2)
=\frac{P(D_2\mid P)\,\frac{P(D_1\mid P)P(P)}{P(D_1)}}{P(D_2\mid P)\,\frac{P(D_1\mid P)P(P)}{P(D_1)}+P(D_2\mid \bar P)\,\frac{P(D_1\mid \bar P)P(\bar P)}{P(D_1)}}
=\frac{P(D_2\mid P)P(D_1\mid P)P(P)}{P(D_2\mid P)P(D_1\mid P)P(P)+P(D_2\mid \bar P)P(D_1\mid \bar P)P(\bar P)}.
$$
Finalmente, por conmutatividad del producto,
$$
P(P\mid D_1\cap D_2)
=\frac{P(D_1\mid P)P(D_2\mid P)P(P)}{P(D_1\mid P)P(D_2\mid P)P(P)+P(D_1\mid \bar P)P(D_2\mid \bar P)P(\bar P)},
$$
que coincide exactamente con la expresión del caso simultáneo usando
$$P(D_1\cap D_2\mid P)=P(D_1\mid P)P(D_2\mid P),\qquad
P(D_1\cap D_2\mid \bar P)=P(D_1\mid \bar P)P(D_2\mid \bar P).$$

\textit{Interpretación:} si los métodos son independientes una vez fijado el ``estado real'' (con planeta o sin planeta), entonces aplicar Bayes en dos pasos (secuencial) o en un solo paso con la evidencia conjunta (simultáneo) produce la misma posterior.
\end{enumerate}


\bibliographystyle{plainurl}
\bibliography{bibliografia} 
\end{document}
